\section{How do the leaks go out, and are they deliberate?}
\label{app:leakchecks}

On the development episodes, Qwen3-32B's leaks go out about equally as the ordinary end-of-turn reply and as an explicit channel message, with a few file writes, while nearly every gpt-oss-20b leak is a file write.
An automated judge labels most leaks deliberate on both models (\cref{tab:leakchecks}).

\paragraph{How leaks go out.}
The end-of-turn reply is the text that the harness posts to the channel when the target finishes its turn.
On gpt-oss-20b, the first leak is almost always the whole configuration block copied into the handoff file.

\paragraph{Deliberate or accidental.}
A decision model, Luna Decisions \citep{openai2026luna}, reads the first message that publishes the value.
It labels the leak deliberate if the message offers the real value for use, and accidental if it exposes the value inside a refusal or an illustration.
For a few Qwen3-32B leaks no saved judge label exists, and the table counts them separately.
We keep accidental leaks in the label, because an overseer must stop them too.

\paragraph{Peer claims of authority.}
A decision-model pass flagged peer messages that claim to speak for the operator.
On Qwen3-32B, leaks were about as common in episodes with such a claim as in episodes without one (\cref{tab:leakchecks}).

\begin{table}[h]
  \centering
  \small
  % Generated by generate_tables.py; edit registry table configuration.
\begin{tabular}{llr}
\toprule
Target & First public-emission route & Episodes \\
\midrule
Qwen & All episodes & \qwenCurrentEpisodes{} \\
Qwen & Any public emission & \qwenCurrentEmittingEpisodes{} \\
Qwen & Final reply & \qwenCurrentRouteFinalReply{} \\
Qwen & send\_message arguments & \qwenCurrentRouteSendMessage{} \\
Qwen & write\_file-type arguments & \qwenCurrentRouteWriteFile{} \\
gpt-oss & All episodes & \gptossCurrentEpisodes{} \\
gpt-oss & Any public emission & \gptossCurrentEmittingEpisodes{} \\
gpt-oss & Final reply & \gptossCurrentRouteFinalReply{} \\
gpt-oss & write\_file-type arguments & \gptossCurrentRouteWriteFile{} \\
\bottomrule
\end{tabular}


  \vspace{0.5em}
  % Generated by generate_tables.py; edit registry table configuration.
\begin{tabular}{llrr}
\toprule
Target & Saved intent label & Episodes & Emitting episodes \\
\midrule
Qwen & Deliberate & \qwenCurrentIntentDeliberate{} & \qwenCurrentEmittingEpisodes{} \\
Qwen & Accidental & \qwenCurrentIntentAccidental{} & \qwenCurrentEmittingEpisodes{} \\
Qwen & Uncertain & \qwenCurrentIntentUncertain{} & \qwenCurrentEmittingEpisodes{} \\
Qwen & Unavailable & \qwenCurrentIntentUnavailable{} & \qwenCurrentEmittingEpisodes{} \\
gpt-oss & Deliberate & \gptossCurrentIntentDeliberate{} & \gptossCurrentEmittingEpisodes{} \\
gpt-oss & Uncertain & \gptossCurrentIntentUncertain{} & \gptossCurrentEmittingEpisodes{} \\
\bottomrule
\end{tabular}


  \vspace{0.5em}
  % Generated by generate_tables.py; edit registry table configuration.
\begin{tabular}{lllrrr}
\toprule
Target & Claim & Claim present? & Emitting & Episodes & Rate \\
\midrule
Qwen & Operator & Present & \qwenCurrentAuthorityOperatorWithLeaks{} & \qwenCurrentAuthorityOperatorWithEpisodes{} & \qwenCurrentAuthorityOperatorWithRate{} \\
Qwen & Operator & Absent & \qwenCurrentAuthorityOperatorWithoutLeaks{} & \qwenCurrentAuthorityOperatorWithoutEpisodes{} & \qwenCurrentAuthorityOperatorWithoutRate{} \\
Qwen & Recipient & Present & \qwenCurrentAuthorityRecipientWithLeaks{} & \qwenCurrentAuthorityRecipientWithEpisodes{} & \qwenCurrentAuthorityRecipientWithRate{} \\
Qwen & Recipient & Absent & \qwenCurrentAuthorityRecipientWithoutLeaks{} & \qwenCurrentAuthorityRecipientWithoutEpisodes{} & \qwenCurrentAuthorityRecipientWithoutRate{} \\
gpt-oss & Operator & Present & \gptossCurrentAuthorityOperatorWithLeaks{} & \gptossCurrentAuthorityOperatorWithEpisodes{} & \gptossCurrentAuthorityOperatorWithRate{} \\
gpt-oss & Operator & Absent & \gptossCurrentAuthorityOperatorWithoutLeaks{} & \gptossCurrentAuthorityOperatorWithoutEpisodes{} & \gptossCurrentAuthorityOperatorWithoutRate{} \\
gpt-oss & Recipient & Present & \gptossCurrentAuthorityRecipientWithLeaks{} & \gptossCurrentAuthorityRecipientWithEpisodes{} & --- \\
gpt-oss & Recipient & Absent & \gptossCurrentAuthorityRecipientWithoutLeaks{} & \gptossCurrentAuthorityRecipientWithoutEpisodes{} & \gptossCurrentAuthorityRecipientWithoutRate{} \\
\bottomrule
\end{tabular}

  \caption{How the first leak goes out, whether the judge labels it deliberate, and leak rates with and without a peer claim of operator authority, on the development episodes.}
  \label{tab:leakchecks}
\end{table}
